\documentclass[11pt, a4paper]{article}
\usepackage{amsmath}
\usepackage{amssymb}
\usepackage{geometry}
\usepackage{accessibility}
\usepackage{booktabs}

% Geometry settings
\geometry{top=2.5cm, bottom=2.5cm, left=2.5cm, right=2.5cm}

\title{\textbf{Verified Mathematical Framework for mmWave Radiotherapy Setup}}
\author{Analysis of Bressler et al. (Med. Phys. 2024)}
\date{\today}

\begin{document}

\maketitle

\begin{abstract}
    This document provides a verified breakdown of the mathematical models used in \textit{Millimeter wave-based patient setup verification and motion tracking during radiotherapy}. The formulas have been cross-referenced with standard radar theory to correct typographical errors found in the original manuscript's text extraction.
\end{abstract}

\section{Core Radar Principles}

\subsection{FMCW Range Estimation}
The fundamental operation of Frequency-Modulated Continuous Wave (FMCW) radar relies on mixing the transmitted ``chirp'' with the received reflection. The resulting Intermediate Frequency (IF) signal contains a beat frequency $\hat{f}_{FFT}$ proportional to the target range $R$.

\textbf{Equation (1): Verified}
\begin{equation}
    \hat{f}_{FFT} = \frac{2 B R}{c T}
    \label{eq:range}
\end{equation}

\noindent \textbf{Definitions:}
\begin{itemize}
    \item $B$: Bandwidth of the chirp (4 GHz).
    \item $c$: Speed of light ($3 \times 10^8$ m/s).
    \item $T$: Chirp duration (Observation window).
\end{itemize}

\subsection{Resolution Limit}
The standard Fast Fourier Transform (FFT) is limited by the observation window $T$. The minimum distinguishable frequency separation is $1/T$. In terms of distance, this creates a resolution limit (Range Bin):
\begin{equation}
    \Delta R_{FFT} = \frac{c}{2B}
\end{equation}
\textit{Verification:} For $B = 4$ GHz, $\Delta R = 3.75$ cm. This matches the paper's reported limitations of the standard FFT approach.

\section{The Chirp Z-Transform (CZT)}

To overcome the coarse resolution of the FFT, the authors utilize the Chirp Z-Transform (CZT). This allows the spectral analysis to focus on a specific frequency band (``zooming in'').

\subsection{Correction of Manuscript Equation (2)}
The manuscript's text contains typographical errors in Equation (2), specifically regarding the indices ($x_s$ vs $x_n$) and the frequency stepping term ($K/K$ vs $k/K$). The corrected mathematical formulation is presented below.

\textbf{Equation (2): Corrected \& Verified}
\begin{equation}
    X_{k, CZT} = \sum_{n=0}^{N-1} x_n e^{-i 2\pi n \left( \frac{f_0 + B \frac{k}{K}}{f_s} \right)}
    \label{eq:czt}
\end{equation}

\noindent \textbf{Correction Notes:}
\begin{itemize}
    \item \textbf{Input Signal:} Changed $x_s$ (typo) to $x_n$, representing the $n$-th sample of the time-domain IF signal.
    \item \textbf{Frequency Step:} Changed $B \frac{K}{K}$ (typo) to $B \frac{k}{K}$, where $k$ is the output frequency bin index ($0 \le k < K$).
    \item \textbf{Physics:} The term inside the exponent represents the phase at time sample $n$ for the specific frequency $f_k = f_0 + B(k/K)$.
\end{itemize}

\subsection{Standard DFT Comparison}
For context, the standard Discrete Fourier Transform (Equation 3 in the paper) is provided below. This scans the entire unit circle rather than a specific arc.
\begin{equation}
    X_{k, DFT} = \sum_{n=0}^{N-1} x_n e^{-i \frac{2\pi n k}{N}}
\end{equation}

\section{Phase-Based Motion Tracking}

While CZT provides high-precision absolute distance, relative motion (displacement) is tracked using the phase of the signal, which is far more sensitive.

\subsection{Phase Extraction}
The phase $\phi$ is extracted directly from the complex output of the CZT at the peak frequency index.

\textbf{Equation (4): Verified}
\begin{equation}
    \phi = \angle \max \left( X_{k, CZT} \right)
\end{equation}

\subsection{Displacement Calculation}
Small displacements cause phase shifts in the reflected signal. The relationship between phase change $\Delta \phi$ and physical displacement $d$ is derived from the Doppler principle.

\textbf{Equation (5): Verified}
\begin{equation}
    d = \frac{c \cdot \Delta \phi}{4\pi f_{min}} = \frac{\lambda_{min} \cdot \Delta \phi}{4\pi}
\end{equation}

\noindent \textbf{Verification of Constants:}
\begin{itemize}
    \item $f_{min}$: The starting frequency of the chirp (77 GHz).
    \item $4\pi$: This factor correctly accounts for the round-trip path of the radar signal. A displacement of $\lambda/2$ results in a round-trip path change of $\lambda$, which corresponds to a full $2\pi$ phase cycle.
\end{itemize}

\section{Summary of Validated Algorithm}
\begin{enumerate}
    \item \textbf{Coarse Search:} Perform FFT to find the approximate range peak frequency $f_{FFT}$.
    \item \textbf{Fine Search (Zoom 1):} Perform CZT centered at $f_{FFT}$ with bandwidth $2F_s/K$.
    \item \textbf{Fine Search (Zoom 2):} Perform CZT centered at the new peak with bandwidth $2F_s/K^2$.
    \item \textbf{Displacement:} Extract phase from the final CZT peak and unwrap it over time to track motion $d$.
\end{enumerate}

\end{document}
